% =====================================================================
% 3. Introduction and Problem Statement
% =====================================================================
\section{Introduction and Problem Statement}
\label{sec:intro}

Battery management systems monitor cell voltages, currents and temperatures, estimate state of charge, balance cells and protect packs from unsafe operating conditions~\cite{plett2015battery}.
Because BMS failures can lead to thermal events in electric vehicles and energy-storage systems, their verification is safety-relevant and is carried out under functional-safety regimes such as ISO~26262~\cite{iso26262}.
Verification is typically performed on hardware-in-the-loop rigs that emulate a battery pack, its loads and its fault conditions while the BMS under test responds in real time.

At the industry partner, the System Test Team (STT) operates such a rig.
Its outputs are currently available only through raw data logs and engineering interfaces, so engineers rely on scripts and offline analysis to interpret results.
This causes two problems identified in the partner's project brief.
First, \textbf{limited visibility}: anomalies during long validation runs may only be found after a test has finished, which lengthens debugging.
Second, \textbf{poor presentation for stakeholders}: there is no interface suitable for showing test status and key performance indicators (KPIs) to project leads or to customers during visits.

Three groups are affected: test engineers, who need detailed, trustworthy real-time data; project leads, who need to see quickly whether a test met its objectives; and customers, who need a clear, credible view of the team's testing rigour.

Building a dashboard is not technically novel in itself.
What is not well understood is how to build one for this setting that is non-intrusive, faithful to noisy signals and appropriate for several audiences at once, and how to show that it is.
Principles for supervisory displays and alarm management are well established for aviation and process control~\cite{sheridan1992telerobotics,isa2016alarm}, but they assume an operator controlling a physical process.
Verification engineering differs: the operator supervises a \emph{test}, judging whether behaviour meets a specification.
Whether those principles transfer to this setting, and how the trade-off between smoothing noise and preserving fault visibility should be handled, has received little attention.
This project addresses that gap with a working artefact used as the vehicle for evaluation.
